\subsection{Symbolic Execution with Refined Query Data}
\label{sec:compare}
After refining query data through the $\sample$, $\extract$, and $\generalize$ stages, the refined query data are reused during symbolic execution of subsequent software releases. This section describes how the $\symb$ procedure (Algorithm~\ref{alg:symbexec}) exploits the refined query data to reduce SMT solving. More specifically, when symbolic execution starts on the next release, $\prevdata$ is provided to the $\Solve$ function (lines~7--8 in Algorithm~\ref{alg:symbexec}). Whenever a new path condition is generated, $\Solve$ first applies the same abstraction introduced in the $\generalize$ stage. For example, suppose $\Solve$ encounters the following path condition:
\[
\pc_{\text{new}}
=
\myset{
(z>28)\land(z<42)\land(z<37)
},
\]
$\Solve$ transforms it into the generalized representation
\[
\pc'_{\text{new}}
=
\myset{
(\avariable_1>\aconstvalue_1)
\land
(\avariable_1<\aconstvalue_1+14)
\land
(\avariable_1<\aconstvalue_1+9)
}.
\]

To efficiently reuse refined query data during symbolic execution, $\Solve$ prioritizes the stored query data that are most likely to match the generalized path condition. Rather than comparing against every stored query data in an arbitrary order, $\Solve$ ranks the stored query data by their similarity to the generalized path condition. The similarity between a generalized path condition and a stored query data is defined as the fraction of generalized constraints in the path condition that are matched by the stored query data. Query data with higher similarity are examined first, increasing the likelihood of finding a matching query data with fewer matching attempts.

Unlike the naive approach, which requires two path conditions to be identical, $\ourtool$ introduces a customized matching method based on feasible-range containment. Specifically, if the feasible range represented by a stored query data is contained within that of the generalized path condition, the stored query data represents a stronger constraint and can be safely reused to determine satisfiability.

For example, consider the generalized path condition
\[
\pc'_{\text{new}}
=
\myset{
(\avariable_1>\aconstvalue_1)
\land
(\avariable_1<\aconstvalue_1+14)
\land
(\avariable_1<\aconstvalue_1+9)
}.
\]
Suppose the stored generalized query data is
\[
\adata
=
(
\myset{
(\avariable_1>\aconstvalue_1)
\land
(\avariable_1<\aconstvalue_1+7)
},
\myset{
\avariable_1\mapsto\aconstvalue_1+1
}
).
\]
Among the three generalized constraints in $\pc'_{\text{new}}$, the constraint
$(\avariable_1>\aconstvalue_1)$
matches exactly, while
the stored constraint
$(\avariable_1<\aconstvalue_1+7)$
is contained in both
$(\avariable_1<\aconstvalue_1+14)$
and
$(\avariable_1<\aconstvalue_1+9)$.
Therefore, all three constraints in $\pc'_{\text{new}}$ are covered by the stored query data, yielding a similarity score of
$\frac{3}{3}=1.0$. Since this entry achieves the highest similarity score, $\Solve$ examines it first. The feasible range represented by $\adata$ is fully contained in that of $\pc'_{\text{new}}$, allowing $\Solve$ to reuse the stored query data. Because the minimum value of $z$ is $28$, $\Solve$ substitutes $\aconstvalue_1=28$ into the generalized assignment, reconstructing the concrete assignment
$\myset{z\mapsto29}$ and immediately returns
$\langle
\texttt{true},
\myset{z\mapsto29}
\rangle$, thereby avoiding an additional SMT solver invocation.

\iffalse
\subsection{Compare Stage} \label{sec:compare}

This stage is designed to effectively match newly explored path conditions in the target release against data from previous versions. The data $\prevdata$, refined through the three stages---$\generalize$, $\extract$, and $\sample$---is provided as input when running symbolic execution on the next release. Specifically, the $\Solve$ function (lines 7--8 in Algorithm~\ref{alg:symbexec}) uses the set $\prevdata$ to effectively solve the path conditions explored in the target release. When a new path condition is provided as input to the $\Solve$ function, $\ourtool$ first applies the $\Generalize$ function to it as those used for query data. For example, consider the following path condition explored in the new release:
{
\[
\scalebox{0.9}{$
\pc_{new} = \myset{(z>28) \land (z<42) \land (z<37)}.
$}
\]
}$\ourtool$ transforms the path condition as follows:
{
\[
\scalebox{0.9}{$
\pc'_{new} = \myset{(\avariable_{1}>\aconstvalue_{1}) \land (\avariable_{1}<\aconstvalue_{1}+14) \land (\avariable_{i}<\aconstvalue_{1}+9)}.
$}
\]
}

With the generalized path condition, we determine a match between the path condition and a query data based on their value ranges and return the result without invoking the solver when a match is found. Recall that the $\Bound$ function derives a tuple $(\avariable, \lowerbound, \upperbound, \nouse)$. To compare a path condition with a query data, we compute the value range of each path condition using the $\Bound$ function. For example, when comparing $\pc'_{new}$ with the following query data:
{
\[
\scalebox{0.9}{$
\adata = (\myset{(\avariable_{1} > \aconstvalue_{1}) \land (\avariable_{1} < \aconstvalue_{1}+7))}, \myset{\avariable_{1} \mapsto \aconstvalue_{1}+1}),
$}
\]
}the $\Bound$ function returns the tuples for the two path conditions as:
{
\[
\scalebox{0.9}{$
\Bound(\pc'_{new}) = (\avariable_{1}, \aconstvalue_{1}, \aconstvalue_{1}+9, \emptyset), \Bound(\pc_{\adata}) = (\avariable_{1}, \aconstvalue_{1}, \aconstvalue_{1}+7, \emptyset).
$}
\]
}If a tuple derived from data entries has a range contained in the new query’s range, indicating that the data entry imposes a stronger condition than the new query, $\ourtool$ considers the query to match those entries and directly returns the satisfiability result. In the example, $\pc'_{new}$ contains the range of $S^{\pc}_{\adata}$, so $\ourtool$ determines that the two path conditions match. It then substitutes $\aconstvalue_{1}=28$ into $\assignments_{\adata}$. Finally, the $\Solve$ function returns $\langle true, \myset{z \mapsto 29} \rangle$. In our experiments, the $\compare$ stage matched 41.5\% more queries on average than the naive approach, which matches a query only when its path condition is identical to that of a data entry.
\fi

